\begin{unipdcode}[language=Python]
class CodaCircolare:
    def __init__(self):
        self._dati = [None] * 4
        self._inizio = 0
        self._dimensione = 0

    def is_empty(self):
        return self._dimensione == 0

    def size(self):
        return self._dimensione

    def enqueue(self, elemento):
        capacita = len(self._dati)
        if self._dimensione == capacita:
            nuovo = [None] * (2 * capacita)
            for i in range(self._dimensione):
                nuovo[i] = self._dati[(self._inizio + i) % capacita]
            self._dati = nuovo
            self._inizio = 0
            capacita *= 2
        indice = (self._inizio + self._dimensione) % capacita
        self._dati[indice] = elemento
        self._dimensione += 1

    def first(self):
        if self.is_empty():
            raise IndexError("Coda vuota")
        return self._dati[self._inizio]

    def dequeue(self):
        elemento = self.first()
        self._dati[self._inizio] = None
        self._inizio = (self._inizio + 1) % len(self._dati)
        self._dimensione -= 1
        return elemento
\end{unipdcode}
